% !TeX root = ../tfg.tex
% !TeX encoding = utf8


\iffalse
################# ASUNCIONES CONDITIONAL MEAN INDEPENDENCE
1. Introducción a la inferencia causal, las asunciones de conditional mean independence,
positivity y exchangeability o SUTVA, y las potential outcomes y el estimado de la diferencia de riesgos y del odds ratio (OR).
positivity
exchangeability o SUTVA
potential outcomes
estimado de la diferencia de riesgos y del odds ratio (OR)
\fi


\subsection{ATE}
\begin{observacion}[La esperanza condicional]
    $\mu_a(x) = E[Y| X = x, A = a]$ es el resultado medio que los pacientes con las covariantes $x$ que fueron
    tratados con el tratamiento $a$.
\end{observacion}

\begin{definicion}[ATE(estadístico)] \label{def:ate}
	Consideramos dos variables aleatorias $X$ e $Y$ y una variable aleatoria binaria $A \in \{0, 1\}$.
	Definimos el \textit{Average Treatment Effect} como
	$$ \psi = E[ E[Y | X, A = 1] - E[Y | X, A = 0 ] ] $$
	Llamaremos a $X$ las covariantes, $A$ el tratamiento e $Y$ el resultado (outcome, en inglés).
	También se puede escribir como
	$$ \psi = E[\mu_1(x) - \mu_0(x)] $$

	%TODO: check "covariantes" is a proper translation, and outcome thing
\end{definicion}
% src: https://alejandroschuler.github.io/mci/7ae35a9031aa4b32b9e82cea3c32a96c.html
\begin{observacion}[Interpretación]
    - El ATE mide la diferencia media entre los pacientes tratados y no tratados, sujeto a las covariantes dadas $x$.
    - El ATE depende de la distribución conjunta de $(Y, X, A)$. Eso es, el ATE depende de la distribución de covariantes $X$.
\end{observacion}

\begin{definicion} Independencia condicional
	Sea $A$ una hipótesis, $B$ y $C$ observaciones.

	Se dice que $A$ y $B$ son condicionalmente independientes dado $C$, y se denota
	$$ (A \perp B )| C $$, si
	$$ P(A | B, C ) = P(A | C)$$
\end{definicion}


Condiciones necesarias para la equivalencia entre el ATE estadístico y ATE causal.
% ( https://onlinelibrary.wiley.com/action/downloadSupplement?doi=10.1002%2Fsim.9234&file=sim9234-sup-0001-supinfo.pdf)
\begin{enumerate}
	\item Consistencia contrafactual.
		$$ Y = A Y(1) + (1 - A)Y(0) $$
		Representa que el resultado observado coincide con el resultado en el modelo contrafactual. Eso es,
		no existe elemento $i$ de la muestra tal que $A_i = 1$ y $Y_i \neq Y(1)$ ni tal quel $A_i = 0$ y $Y_i \neq Y(0)$.
	\item Adicionalmente a lo anterior, asumimos observaciones independientes y no error en la medida. % TODO: fix wording
	\item Intercambiabilidad condicional, independencia condicional o ignorabilidad % TODO; check understanding
		Independencia entre $Y(a)$ y $A$ bajo la distribución condicionada a $C$ % TODO: check is correct interpretation (associatively)
		No confundir con indepedencia entre $Y(a)$ y $A | C$
	\item Positividad
		$$ P(A | C = c) > 0 \; \forall c \text{tales que}  P(C = c) > 0$$
\end{enumerate}

Bajo las condiciones anteriores, por el \ref{thm:probabilidad-total} tenemos que
$$ P[Y(a) = 1 ] = \sum_c P[Y(a) = 1 | C = c]P[C=c] $$
Aplicamos la intercambiabilidad condicional, por la cual dentro de cada estrato $c$ de $C$ los predictores están
equitativamente distribuidos entre el grupo tratado ($A = 1$)y no tratado ($A = 0$):
$$ = \sum_c P[Y(a) = 1 | A=a,C=c]P[C=c] $$
